\newpage
\chapter{KESIMPULAN DAN SARAN} \label{Bab V}

\section{Kesimpulan} \label{V.Kesimpulan}
Berdasarkan hasil penelitian yang telah dilakukan, dapat disimpulkan bahwa penelitian ini berhasil menjawab seluruh rumusan masalah melalui perancangan, implementasi, serta evaluasi model klasifikasi berbasis multimodal. Proses penelitian dimulai dari pengumpulan dan anotasi data, pembangunan model unimodal dan multimodal, hingga analisis performa menggunakan berbagai skenario eksperimen. Hasil yang diperoleh menunjukkan bahwa pendekatan multimodal mampu memberikan peningkatan performa yang signifikan dibandingkan pendekatan unimodal. Selain itu, penelitian ini juga memberikan \textit{insight} terkait efektivitas metode \textit{fusion} serta kompleksitas arsitektur dalam pemodelan hubungan antar modalitas. Adapun kesimpulan yang diperoleh adalah sebagai berikut:

\begin{enumerate}
    \item Penelitian ini berhasil membangun model klasifikasi biner berbasis multimodal dengan mengintegrasikan CLIP sebagai \textit{encoder} visual dan ELECTRA sebagai \textit{encoder} tekstual melalui pendekatan \textit{intermediate fusion}. Model yang dikembangkan mampu memproses dan menggabungkan informasi visual dan tekstual secara simultan, sehingga dapat memahami konteks meme secara lebih komprehensif. Dengan memanfaatkan kekuatan masing-masing \textit{encoder}, model mampu menangkap pola visual dari gambar serta makna linguistik dari teks yang saling melengkapi dalam proses klasifikasi.
    
    \item Integrasi representasi visual dan tekstual melalui pendekatan multimodal terbukti memberikan peningkatan performa yang signifikan dibandingkan model unimodal. Hal ini ditunjukkan oleh peningkatan nilai \textit{F1-score} dari 91.5\% pada model unimodal gambar dan 92.0\% pada model unimodal teks menjadi 96.1\% pada model multimodal. Peningkatan ini menunjukkan bahwa kombinasi kedua modalitas mampu memberikan pemahaman konteks yang lebih baik, terutama pada data meme yang mengandung makna implisit. Dengan demikian, pendekatan multimodal lebih efektif dalam menangani permasalahan klasifikasi yang melibatkan hubungan kompleks antara gambar dan teks.
    
    \item Evaluasi kinerja model menggunakan \textit{confusion matrix} menunjukkan bahwa model multimodal memiliki performa yang baik dalam mendeteksi konten \textit{self-harm}, khususnya dalam meningkatkan kemampuan deteksi pada kelas minoritas. Model menghasilkan nilai \textit{precision} yang tinggi serta \textit{recall} yang lebih baik dibandingkan model unimodal, yang menunjukkan bahwa model mampu mengurangi kesalahan klasifikasi, terutama \textit{false negative} yang penting dalam konteks deteksi \textit{self-harm}. Selain itu, hasil eksperimen juga menunjukkan bahwa penggunaan arsitektur yang lebih kompleks seperti \textit{Transformer Encoder} tidak memberikan peningkatan performa yang signifikan dibandingkan metode \textit{simple fusion}. Hal ini mengindikasikan bahwa pendekatan \textit{simple fusion} sudah cukup efektif dalam menangkap hubungan antar modalitas, sehingga peningkatan kompleksitas model tidak selalu diperlukan dalam konteks permasalahan ini.
\end{enumerate}


\section{Saran} \label{V.Saran}
Berdasarkan penelitian yang telah dilakukan, maka saran yang dapat diajukan adalah sebagai berikut:


\begin{enumerate}
    \item Penelitian selanjutnya disarankan untuk melibatkan ahli di bidang psikologi atau kesehatan mental dalam proses verifikasi hasil anotasi data. Meskipun pendekatan \textit{ensemble pseudo-labeling} yang digunakan dalam penelitian ini mampu menghasilkan anotasi dengan tingkat kesepakatan yang tinggi, validasi oleh ahli tetap diperlukan untuk memastikan akurasi dan sensitivitas dalam mengidentifikasi konten \textit{self-harm}, terutama pada kasus yang bersifat implisit atau ambigu. Keterlibatan ahli psikologi diharapkan dapat meningkatkan kualitas label dataset, sehingga model yang dihasilkan menjadi lebih reliabel dan sesuai dengan konteks klinis maupun sosial yang sebenarnya.
    
    
 
    
    \item Selain itu, penelitian selanjutnya dapat mengembangkan sistem deteksi yang lebih aplikatif, seperti integrasi model ke dalam sistem monitoring media sosial secara \textit{real-time}. Implementasi ini dapat dimanfaatkan sebagai langkah preventif dalam mengidentifikasi dan menangani konten \textit{self-harm} secara lebih dini.
\end{enumerate}